% GENERATED FILE. Edit canonical lessons, then run npm run generate:books.
% canonical-source-hash: fnv1a64:f9bed786bc3f9acd
% canonical-lessons: GE-C72-abendessen, GE-C72-fruehstueck, GE-C72-reis, GE-C72-birne, GE-C72-zitrone

\chapter{Dinner, Breakfast, Rice, Pear, Lemon}
\label{ch:ge-dinner-breakfast-rice-pear-lemon}
\hlchaptermodality{72}

\begin{chapteropening}
\textbf{By the end of this chapter:} \emph{I can say dinner, breakfast, rice, a pear and a lemon.}

Complete the last lesson of chapter 72: I can say dinner, breakfast, rice, a pear and a lemon.
\end{chapteropening}

\section[das Abendessen]{das Abendessen — dinner}

\label{lesson:GE-C72-abendessen}



\begin{quote}
Before the new one: say the German for oil, then the German for lunch.
\end{quote}

\subsection*{You'll want to know: das Abendessen}
\textbf{das Abendessen} — \textquotedblleft{}dinner\textquotedblright{}.

\begin{cousinweb}
\textbf{Abendessen} is \emph{Abend} + the word for \textquotedblleft{}a meal\textquotedblright{}, which is the verb \textquotedblleft{}to eat\textquotedblright{} $\leftarrow$ Old High German \emph{ezzan} $\leftarrow$ Proto-Germanic \emph{*etaną}, English \emph{eat}. (EtymDB 2.1)
\end{cousinweb}

\begin{grammarlens}[title={What you've built}]
One more word for this chapter.
\end{grammarlens}

\subsection*{Your turn}
\begin{itemize}
\raggedright
  \item \emph{Say it:} \emph{das Abendessen}
  \item \emph{Say it:} \emph{das Abendessen}, once more
  \item \emph{Say it:} say \emph{das Abendessen}
  \item \emph{Recall:} say the German for oil, then the German for lunch, then say \emph{das Abendessen} again
\end{itemize}

\begin{tcolorbox}[breakable,colback=teal!4,colframe=teal!35!black,title={Before you move on}]
What does das Abendessen mean? (\textquotedblleft{}Dinner\textquotedblright{}.) Say it once more.
\end{tcolorbox}
\section[das Frühstück]{das Frühstück — breakfast}

\label{lesson:GE-C72-fruehstueck}



\begin{quote}
Before the new one: say the German for lunch, then the German for dinner.
\end{quote}

\subsection*{You'll want to know: das Frühstück}
\textbf{das Frühstück} — \textquotedblleft{}breakfast\textquotedblright{}.

\begin{cousinweb}
\textbf{Frühstück} is \emph{früh}, \textquotedblleft{}early\textquotedblright{}, + the word for \textquotedblleft{}a piece\textquotedblright{}, Old High German \emph{stucki} (EtymDB 2.1): the early bite of bread.
\end{cousinweb}

\begin{grammarlens}[title={What you've built}]
One more word for this chapter.
\end{grammarlens}

\subsection*{Your turn}
\begin{itemize}
\raggedright
  \item \emph{Say it:} \emph{das Frühstück}
  \item \emph{Say it:} \emph{das Frühstück}, once more
  \item \emph{Say it:} say \emph{das Frühstück}
  \item \emph{Recall:} say the German for lunch, then the German for dinner, then say \emph{das Frühstück} again
\end{itemize}

\begin{tcolorbox}[breakable,colback=teal!4,colframe=teal!35!black,title={Before you move on}]
What does das Frühstück mean? (\textquotedblleft{}Breakfast\textquotedblright{}.) Say it once more.
\end{tcolorbox}
\section[der Reis]{der Reis — rice}

\label{lesson:GE-C72-reis}



\begin{quote}
Before the new one: say the German for dinner, then the German for breakfast.
\end{quote}

\subsection*{You'll want to know: der Reis}
\textbf{der Reis} — \textquotedblleft{}rice\textquotedblright{}.

\begin{cousinweb}
\textbf{Reis} $\leftarrow$ Middle High German \emph{rīs} $\leftarrow$ Latin \emph{oryza} $\leftarrow$ Greek \emph{óryza} — borrowed, like English \emph{rice}. (EtymDB 2.1 has the Latin and Greek step; its German \emph{Reis} row is a different word, \textquotedblleft{}twig\textquotedblright{}.)
\end{cousinweb}

\begin{grammarlens}[title={What you've built}]
One more word for this chapter.
\end{grammarlens}

\subsection*{Your turn}
\begin{itemize}
\raggedright
  \item \emph{Say it:} \emph{der Reis}
  \item \emph{Say it:} \emph{der Reis}, once more
  \item \emph{Say it:} say \emph{der Reis}
  \item \emph{Recall:} say the German for dinner, then the German for breakfast, then say \emph{der Reis} again
\end{itemize}

\begin{tcolorbox}[breakable,colback=teal!4,colframe=teal!35!black,title={Before you move on}]
What does der Reis mean? (\textquotedblleft{}Rice\textquotedblright{}.) Say it once more.
\end{tcolorbox}
\section[die Birne]{die Birne — a pear}

\label{lesson:GE-C72-birne}



\begin{quote}
Before the new one: say the German for breakfast, then the German for rice.
\end{quote}

\subsection*{You'll want to know: die Birne}
\textbf{die Birne} — \textquotedblleft{}a pear\textquotedblright{}.

\begin{cousinweb}
\textbf{Birne} $\leftarrow$ Middle High German \emph{bir} (EtymDB 2.1), from Latin \emph{pira}, \textquotedblleft{}pears\textquotedblright{} — borrowed, like English \emph{pear}.
\end{cousinweb}

\begin{grammarlens}[title={What you've built}]
One more word for this chapter.
\end{grammarlens}

\subsection*{Your turn}
\begin{itemize}
\raggedright
  \item \emph{Say it:} \emph{die Birne}
  \item \emph{Say it:} \emph{die Birne}, once more
  \item \emph{Say it:} say \emph{die Birne}
  \item \emph{Recall:} say the German for breakfast, then the German for rice, then say \emph{die Birne} again
\end{itemize}

\begin{tcolorbox}[breakable,colback=teal!4,colframe=teal!35!black,title={Before you move on}]
What does die Birne mean? (\textquotedblleft{}A pear\textquotedblright{}.) Say it once more.
\end{tcolorbox}
\section[die Zitrone]{die Zitrone — a lemon}

\label{lesson:GE-C72-zitrone}



\begin{quote}
Before the new one: say the German for rice, then the German for a pear.
\end{quote}

\subsection*{You'll want to know: die Zitrone}
\textbf{die Zitrone} — \textquotedblleft{}a lemon\textquotedblright{}.

\begin{cousinweb}
\textbf{Zitrone} $\leftarrow$ French \emph{citron}, \textquotedblleft{}lemon\textquotedblright{} $\leftarrow$ Latin \emph{citrus} (EtymDB 2.1) — borrowed.
\end{cousinweb}

\begin{grammarlens}[title={What you've built}]
One more word for this chapter.
\end{grammarlens}

\subsection*{Your turn}
\begin{itemize}
\raggedright
  \item \emph{Say it:} \emph{die Zitrone}
  \item \emph{Say it:} \emph{die Zitrone}, once more
  \item \emph{Say it:} say \emph{die Zitrone}
  \item \emph{Recall:} say the German for rice, then the German for a pear, then say \emph{die Zitrone} again
\end{itemize}

\begin{tcolorbox}[breakable,colback=teal!4,colframe=teal!35!black,title={Before you move on}]
What does die Zitrone mean? (\textquotedblleft{}A lemon\textquotedblright{}.) Say it once more.
\end{tcolorbox}
